\newpage
\chapter*{CONCLUSION GÉNÉRALE}
\addcontentsline{toc}{chapter}{CONCLUSION GÉNÉRALE}

Ce travail de recherche a été consacré à l'optimisation des transmissions \textbf{LR-FHSS} à l'aide de l'apprentissage par renforcement profond (\textit{Deep Reinforcement Learning}). Face aux défis croissants de passage à l'échelle et d'efficacité énergétique des réseaux IoT, nous avons proposé une approche dynamique capable d'adapter intelligemment les paramètres de transmission aux conditions changeantes du canal radio.

La première phase de notre étude a permis la conception et la validation d'un simulateur événementiel développé en Python. En respectant rigoureusement les spécifications techniques de la norme LR-FHSS et en confrontant nos résultats à des données réelles, nous avons obtenu une erreur quadratique moyenne (RMSE) de \textbf{12,14 \%}, garantissant ainsi la fiabilité des évaluations de performances menées par la suite.

L'implémentation d'un agent fondé sur l'algorithme \textbf{Double Deep Q-Network (DDQN)} a ensuite démontré la pertinence de l'intelligence artificielle face aux méthodes d'allocation statiques. Nos résultats mettent en évidence une amélioration du taux de livraison des paquets (PDR) de \textbf{2,48\%} en moyenne, tout en réalisant une réduction significative de la puissance d'émission de \textbf{32 \%}. Cette optimisation se traduit par une diminution de la consommation énergétique globale de \textbf{4,9 \%}, confirmant la viabilité de cette approche pour prolonger l'autonomie des nœuds capteurs dans des environnements contraints.

Malgré ces avancées significatives, ce travail ouvre la voie à plusieurs perspectives de recherche. Une modélisation encore plus fine du canal radio pourrait être envisagée en intégrant des modèles d'évanouissement (\textit{fading}) complexes, ainsi qu'une gestion plus précise des variations temporelles de l'ombrage (\textit{shadowing}). Par ailleurs, l'inclusion du surcoût protocolaire réel et des mécanismes de retransmission permettrait d'affiner le bilan énergétique global du système.

En somme, cette étude démontre que l'apprentissage par renforcement constitue un levier technologique indispensable pour l'évolution des réseaux LPWAN vers plus d'autonomie et d'efficacité spectrale. À terme, le passage de la simulation au déploiement sur des passerelles de production représentera une étape décisive pour soutenir la croissance de l'Internet des Objets.